% --- PREFACE / SCOPE ---
\addchap{Preface: Scope and Purpose}

\noindent \textbf{The Motivation}\\
These notes were initially born out of a personal need: the desire to organize, synthesize, and rigorously typeset the vast amount of information presented during the lectures. Writing and structuring this material in \LaTeX\ served as my primary method of active recall and study. However, recognizing how challenging it can be to study from scattered slides, recordings, and fragmented notes, I decided to open-source this project. The goal is to provide current and future students of the Master's Degree in Computer Science with a unified, structured, and accessible study companion.

\vspace{1em}
\noindent \textbf{The Scope}\\
The scope of this document is to cover the core theoretical concepts, definitions, vulnerabilities, and software life-cycle processes discussed throughout the course. It is designed to act as a comprehensive bridge between the professors' official slides and the verbal explanations provided during class. 

\vspace{1em}
\noindent \textbf{How to Use This Material}\\
While I have made every effort to ensure the technical rigor of this document, it is important to state what these notes are \textit{not}. They are \textbf{not a replacement} for the official recommended textbooks, nor are they a substitute for attending lectures and engaging with the professors. They are a student-made study guide. You should use them to clarify complex topics, review for exams, and quickly reference formal definitions.

\vspace{1em}
\noindent \textbf{Collaborative Spirit}\\
Because this is an actively maintained, open-source project, it is subject to continuous improvement. If you are studying from this document and spot a typographical error, a technical inaccuracy, or a missing concept, please do not hesitate to reach out via the contact information on the title page. Collaborative corrections are highly encouraged and deeply appreciated.

\vspace{2em}
\noindent \textit{Happy studying,}\\
\textbf{Luca Facchini}\\
BSc Graduate in Computer Science
